\chapter{Choking (Airway Obstruction)}

\section*{Definition}
Partial or complete blockage of the airway by a foreign object, food, or anatomical structure, preventing adequate breathing and requiring immediate intervention.

\section*{When to See a Doctor}
See your doctor if:
\begin{itemize}
  \item Complete airway obstruction (unable to cough, speak, or breathe) requires immediate emergency response and choking first aid
\end{itemize}

\section*{How It Develops}
Foreign body lodges in the pharynx, larynx, or trachea causing mechanical obstruction. Partial obstruction allows some air exchange (coughing, wheezing, stridor). Complete obstruction causes silent distress with inability to cough, speak, or breathe, leading to hypoxia, loss of consciousness, and cardiac arrest within minutes without intervention.

\section*{Causes}
\begin{itemize}
  \item Food aspiration (meat, nuts, seeds, hot dogs, grapes)
  \item Small objects (coins, marbles, toy parts, buttons)
  \item Inadequate chewing or eating too fast
  \item Alcohol intoxication (impaired swallowing reflex)
  \item Dental issues or poorly fitting dentures
  \item Neurological conditions (dysphagia from stroke, Parkinson, dementia)
  \item Anaphylaxis (laryngeal edema causing airway closure)
  \item Croup or epiglottitis (infectious airway narrowing)
  \item Trauma causing laryngeal fracture or airway compression
  \item Tongue obstruction in unconscious patient
\end{itemize}

\section*{Related Symptoms}
\begin{itemize}
  \item Shortness of breath
  \item Stridor
  \item Cough
  \item Cyanosis
  \item Loss of consciousness
\end{itemize}

\section*{Medical Evaluation}
\begin{itemize}
  \item After an acute choking episode, your doctor will assess airway patency, oxygen saturation, and examine the oropharynx for residual foreign body or mucosal injury.
  \item Neck and chest X-rays (anteroposterior and lateral) are obtained to identify radiopaque foreign bodies in the airway or oesophagus.
  \item Flexible nasolaryngoscopy or bronchoscopy is performed for direct visualisation and removal of non-radiopaque foreign bodies or those lodged in the larynx or trachea.
  \item Referral to an otolaryngologist or gastroenterologist is arranged for endoscopic removal of foreign bodies not accessible in the emergency setting.
\end{itemize}

\section*{Treatment Options}
\begin{itemize}
  \item Emergency foreign body removal is performed using direct laryngoscopy, bronchoscopy, or oesophagoscopy depending on the location of the obstruction.
  \item Post-obstruction management includes monitoring for aspiration pneumonia with chest X-ray and antibiotics if needed.
  \item Swallowing assessment by a speech-language pathologist is arranged for patients with recurrent choking or known dysphagia.
  \item Dietary modifications (soft or pureed foods, thickened liquids) and swallowing manoeuvres are taught to reduce aspiration risk in patients with neurological dysphagia.
\end{itemize}

\section*{Self-Care and Home Management}
\begin{itemize}
  \item For a conscious adult or child with complete airway obstruction, call emergency services and give cycles of five back blows followed by five abdominal thrusts; use chest thrusts instead of abdominal thrusts for an infant or a pregnant or markedly obese person.
  \item If alone and choking, perform abdominal thrusts on yourself using a chair back or fist, or lean forward over a firm edge.
  \item For partial obstruction where the person can cough or speak, encourage continued coughing and monitor closely; do not perform thrusts while an effective cough is present.
  \item After the object is expelled or if the person becomes unconscious, begin CPR starting with chest compressions and check the mouth for visible objects.
  \item Learn basic life support (BLS) and choking first aid through a certified training course to be prepared for emergencies.
\end{itemize}

\section*{References}
\begin{thebibliography}{99}
\bibitem{choking:ref1} Kline RM. ``Choking and foreign body aspiration.'' In: Nelson Textbook of Pediatrics. 22nd ed. Elsevier; 2024.
\bibitem{choking:ref2} American Heart Association. ``2020 American Heart Association guidelines for CPR and ECC: management of foreign-body airway obstruction.'' Circulation. 2020;142(16 Suppl 2):S366-S468.
\bibitem{choking:ref3} Darrow DH, Holinger LD. ``Foreign bodies of the airway.'' In: Cummings Otolaryngology. 7th ed. Elsevier; 2020.
\bibitem{choking:ref4} Lichtenstein R, Thompson L. ``Foreign body aspiration in adults.'' Am Fam Physician. 2021;104(6):589-595.
\bibitem{choking:ref5} Alikorn R, Chen X, Milkovich S, et al. ``Fatal and non-fatal choking incidents in children.'' Pediatrics. 2019;144(4):e20190388.
\bibitem{choking:ref6} Jaffe RB. ``Airway foreign bodies in children.'' UpToDate. Wolters Kluwer; 2025.
\end{thebibliography}
